% Zone „Das kennst du schon“ – aus bank/quadratische-funktionen/zone.jsonl (zusammenbau v0.1)
\einheitenkopf[zone][Kennst du schon]{Das kennst du schon}
\setcounter{aufgabe}{0}

%% TODO Titel: Ich-kann-Satz fehlt in der Bank (Platzhalter: Kettenname) – Nr. 1
%% TODO Anweisung: ein Satz über der ersten Teilaufgabe fehlt in der Bank – Nr. 1
\begin{aufgabe}{Quadrieren auch negativer Zahlen und Dezimalzahlen, Quadrat vor Punkt vor Strich: (−7)² = 49, nicht −49; 0,5² = 0,25}
\begin{teile}
\teil $6^2$ – wie viel? \leerfeld
\teil $1{,}2^2$ – wie viel? \leerfeld
\end{teile}
\end{aufgabe}

%% TODO Titel: Ich-kann-Satz fehlt in der Bank (Platzhalter: Kettenname) – Nr. 2
%% TODO Anweisung: ein Satz über der ersten Teilaufgabe fehlt in der Bank – Nr. 2
\begin{aufgabe}{Koordinaten lesen und eintragen, vier Quadranten, (x | y)-Reihenfolge, Kästchenraster mit 0,5-Einteilung}
\begin{teile}
\teil Punkt $A$ – Koordinaten?

\begin{ksys}[xmin=-4,xmax=4,ymin=-4,ymax=4,ablesen]
\punkt{3}{2}{A}
\end{ksys}
A( \leerfeld | \leerfeld )
\teil Punkt $C$ – Koordinaten?

\begin{ksys}[xmin=-5,xmax=4,ymin=-4,ymax=4,ablesen]
\punkt{-4}{-3}{C}
\end{ksys}
C( \leerfeld | \leerfeld )
\end{teile}
\end{aufgabe}

%% TODO Titel: Ich-kann-Satz fehlt in der Bank (Platzhalter: Kettenname) – Nr. 3
%% TODO Anweisung: ein Satz über der ersten Teilaufgabe fehlt in der Bank – Nr. 3
\begin{aufgabe}{Lineare Funktion f(x) = m·x + n: Gerade zeichnen, Funktionswert, Punktprobe, Nullstelle}
\begin{teile}
\teil $g(x) = 2x + 1$ – $g(4)$? \leerfeld
\teil $g(x) = -3x + 2$ – $g(-2)$? \leerfeld
\end{teile}
\end{aufgabe}

%% TODO Titel: Ich-kann-Satz fehlt in der Bank (Platzhalter: Kettenname) – Nr. 4
%% TODO Anweisung: ein Satz über der ersten Teilaufgabe fehlt in der Bank – Nr. 4
\begin{aufgabe}{Binomische Formel ausmultiplizieren und zusammenfassen, auch mit Minus in der Klammer}
\begin{teile}
\teil $(x + 4)^2$ – ausmultipliziert? \leerfeld
\teil $(x - 6)^2$ – ausmultipliziert? \leerfeld
\end{teile}
\end{aufgabe}

%% TODO Titel: Ich-kann-Satz fehlt in der Bank (Platzhalter: Kettenname) – Nr. 5
%% TODO Anweisung: ein Satz über der ersten Teilaufgabe fehlt in der Bank – Nr. 5
\begin{aufgabe}{Lineare Gleichung lösen, Terme mit x auf eine Seite bringen, Schreibform mit Strich}
\begin{gleichungsraster}[2]
\gl{\text{$3x + 2 = 14$ – $x$?}} &
\gl{\text{$2x + 9 = -3$ – $x$?}} \\
\end{gleichungsraster}
\end{aufgabe}

%% TODO Titel: Ich-kann-Satz fehlt in der Bank (Platzhalter: Kettenname) – Nr. 6
%% TODO Anweisung: ein Satz über der ersten Teilaufgabe fehlt in der Bank – Nr. 6
\begin{aufgabe}{Schnittpunkt zweier Geraden durch Gleichsetzen}
\begin{gleichungsraster}[2]
\gl{\text{$g(x) = x + 1$ und $h(x) = -x + 5$ – Schnittpunkt?}} &
\gl{\text{$g(x) = 3x + 4$ und $h(x) = x - 2$ – Schnittpunkt?}} \\
\end{gleichungsraster}
\end{aufgabe}

%% TODO Titel: Ich-kann-Satz fehlt in der Bank (Platzhalter: Kettenname) – Nr. 7
%% TODO Anweisung: ein Satz über der ersten Teilaufgabe fehlt in der Bank – Nr. 7
\begin{aufgabe}{Quadratwurzel ziehen, auch mit dem Taschenrechner, Näherungswert runden; aus einer negativen Zahl gibt es keine Wurzel}
\begin{teile}
\teil $\sqrt{64}$ – wie viel? \leerfeld
\teil $\sqrt{7}$ – auf zwei Stellen gerundet? \leerfeld
\end{teile}
\end{aufgabe}

%% TODO Titel: Ich-kann-Satz fehlt in der Bank (Platzhalter: Kettenname) – Nr. 8
%% TODO Anweisung: ein Satz über der ersten Teilaufgabe fehlt in der Bank – Nr. 8
\begin{aufgabe}{Quadratische Gleichung durch Wurzelziehen lösen und Normalform mit der p-q-Formel lösen, beide Lösungen angeben}
\begin{gleichungsraster}[2]
\gl{\text{$x^2 = 36$ – Lösungen?}} &
\gl{\text{$x^2 + 4x - 5 = 0$ – Lösungen?}} \\
\end{gleichungsraster}
\end{aufgabe}

%% TODO Titel: Ich-kann-Satz fehlt in der Bank (Platzhalter: Kettenname) – Nr. 9
%% TODO Anweisung: ein Satz über der ersten Teilaufgabe fehlt in der Bank – Nr. 9
\begin{aufgabe}{Binomische Formel ausmultiplizieren und zusammenfassen, auch mit Minus in der Klammer – Fehler finden}
\begin{teile}
\teil Tim rechnet so: \rechnung{(x - 5)^2 &= x^2 - 25} Finde den Fehler und rechne richtig. \leerfeld
\end{teile}
\end{aufgabe}

%% TODO Titel: Ich-kann-Satz fehlt in der Bank (Platzhalter: Kettenname) – Nr. 10
%% TODO Anweisung: ein Satz über der ersten Teilaufgabe fehlt in der Bank – Nr. 10
\begin{aufgabe}{Binomische Formel ausmultiplizieren und zusammenfassen, auch mit Minus in der Klammer – selbst rechnen}
\begin{teile}
\teil $(x + 7)^2$ – ausmultipliziert? \leerfeld
\end{teile}
\end{aufgabe}
